\chapter{Réseaux transrégionaux des cryptomonnaies}\label{ch:nw:cross-region-crypto-networks}

La liquidité des cryptomonnaies est répartie sur quelques régions cloud, et une société qui négocie à travers elles a besoin que ses messages traversent des océans. En avril 2026,
un aller-retour de Tokyo à Francfort par la dorsale propre du fournisseur cloud, par peering VPC, a pris une médiane de \qty{224.3}{\milli\second}~;
par le Direct Connect du même fournisseur, relié à un réseau spécialisé, \qty{136.5}{\milli\second}. La lumière dans une fibre droite, le long de la
géodésique, mettrait \qty{91.3}{\milli\second}. La dorsale du fournisseur est conçue pour «~la résilience et l'échelle~», dit son blog, et elle prend le
long chemin~; le spécialiste en prend un plus court~; et la société qui ne peut pas choisir envoie le même message par les deux et garde le premier arrivé.

Les chapitres~16 à~18 ont travaillé à l'intérieur d'une seule région. Ce chapitre relie les régions~: les hubs et leurs distances, la dorsale et
l'internet public face aux lignes privées et aux fournisseurs spécialisés, et la \omterm{def:nw:cross-region-crypto-networks:idem}{mise en concurrence des routes} (route racing), avec les ordres en double
qu'elle crée et ce que les règles d'une plateforme en font. Les chiffres publiés sont cités et datés~; la mise en concurrence est une simulation étiquetée.

\section{Les hubs~: Tokyo, Singapour, Hong Kong, Londres, Virginie}

\begin{definition}[Allongement du chemin]\label{def:nw:cross-region-crypto-networks:infl}
L'\emph{allongement du chemin}\index{allongement du chemin} d'une route entre deux lieux est son aller-retour mesuré divisé par l'aller-retour que la lumière
mettrait le long de la géodésique dans le milieu de la route (pour les réseaux longue distance, la fibre)~: c'est le \omterm{def:nw:the-north-american-map:factor}{facteur de tracé} du chapitre~10, pour les
aller-retours sur des réseaux dont le milieu et le tracé ne sont pas publiés.
\end{definition}

Les hubs sont ceux du chapitre~17~: Tokyo, où la région ap-northeast-1 d'AWS héberge Binance, BitMEX et les validateurs de Hyperliquid~; Singapour, la région
de Bybit~; Hong Kong~; Londres et Francfort en Europe~; le nord de la Virginie aux États-Unis. Le fournisseur ne publie que les noms des régions et leur
nombre de zones (quatre à Tokyo, six en Virginie du Nord, trois dans les autres ici), si bien que le chapitre place chacune dans sa ville, et la Virginie du Nord
à Ashburn, une hypothèse signalée dans le tableau des sites.

\begin{table}[ht]
\centering\small
\begin{tabular}{lrrrrrrr}
\toprule
De Tokyo à & Géodésique & Vide & Fibre & Dorsale & Spécialiste & \multicolumn{2}{c}{\omterm{def:nw:cross-region-crypto-networks:infl}{Allongement du chemin}} \\
 & (km) & RTT (ms) & RTT (ms) & (ms) & (ms) & dorsale & spécialiste \\
\midrule
Hong Kong & 2\,883 & 19.2 & 28.1 & -- & -- & -- & -- \\
Singapour & 5\,311 & 35.4 & 51.8 & 68.0 & 65.4 & 1.31 & 1.26 \\
N. Virginie & 10\,895 & 72.7 & 106.3 & 150.2 & 135.4 & 1.41 & 1.27 \\
Francfort & 9\,359 & 62.4 & 91.3 & 224.3 & 136.5 & 2.46 & 1.50 \\
Londres & 9\,586 & 64.0 & 93.5 & 210.0 & 139.6 & 2.25 & 1.49 \\
Stockholm & 8\,194 & 54.7 & 79.9 & 245.3 & 124.3 & 3.07 & 1.56 \\
\bottomrule
\end{tabular}
\caption{Les hubs de Tokyo~: \omterm{def:nw:the-north-american-map:geodesic}{distances géodésiques} entre les villes des régions et les aller-retours que la lumière met le long de celles-ci, face aux aller-retours
médians publiés par AWS pour avril 2026 sur sa dorsale (peering VPC) et sur Direct Connect avec un réseau spécialisé, et l'\omterm{def:nw:cross-region-crypto-networks:infl}{allongement du chemin} de chacun par
rapport au plancher de la fibre. Données~: \texttt{nw\_race.hub\_rows()}.}\label{tab:nw:cross-region-crypto-networks:hubs}
\end{table}

Le motif du tableau est le sujet du chapitre. Vers Singapour et la Virginie, les deux réseaux passent 26 à 41~\% au-dessus de la lumière dans une fibre droite. Vers
l'Europe, la dorsale va à 2.2 à 3.1 fois le plancher, le spécialiste à 1.5 fois. Ni l'un ni l'autre ne publie son tracé~; l'allongement indique que la dorsale fait
un long détour et le spécialiste un plus court. La meilleure route européenne depuis Tokyo, \qty{124.3}{\milli\second} jusqu'à Stockholm, reste plus de deux fois
le plancher du vide.

\begin{omfigure}
\begin{tikzpicture}
\begin{axis}[ybar,width=11.5cm,height=6cm,bar width=8pt,ymin=0,ymax=270,xtick=data,
  symbolic x coords={Singapore,N. Virginia,Frankfurt,London,Stockholm},xticklabels={Singapour,N. Virginie,Francfort,Londres,Stockholm},ylabel={aller-retour depuis Tokyo (ms)},
  tick label style={font=\footnotesize},label style={font=\small},grid=major,grid style={gray!25},legend columns=3,
  legend style={font=\footnotesize,at={(0.5,-0.18)},anchor=north,draw=none},table/col sep=comma,enlarge x limits=0.15]
\addplot[fill=gray!40,draw=gray] table[x=hub,y=fibre_floor]{figdata/networks/19-cross-region-crypto-networks/hubs.csv};
\addplot[fill=omDef!55,draw=omDef] table[x=hub,y=specialist]{figdata/networks/19-cross-region-crypto-networks/hubs.csv};
\addplot[fill=omThm!45,draw=omThm] table[x=hub,y=backbone]{figdata/networks/19-cross-region-crypto-networks/hubs.csv};
\legend{plancher de la fibre (géodésique),réseau spécialisé,dorsale du fournisseur}
\end{axis}
\end{tikzpicture}

\omcaption{Aller-retours depuis Tokyo~: le plancher de la fibre le long de la géodésique, et les médianes publiées sur un réseau spécialisé et sur la dorsale du
fournisseur (avril 2026). Données~: \texttt{fig\_race.py}, la \cref{tab:nw:cross-region-crypto-networks:hubs}.}\label{fig:nw:cross-region-crypto-networks:hubs}
\end{omfigure}

\section{Dorsale cloud, internet public et lignes privées}

\begin{definition}[Border Gateway Protocol, dorsale cloud]\label{def:nw:cross-region-crypto-networks:bgp}
Le \emph{Border Gateway Protocol}\index{Border Gateway Protocol} (BGP) est le protocole de routage de l'internet entre \omterm{def:nw:where-the-crypto-venues-live:as}{systèmes autonomes}~: chacun échange des informations
d'accessibilité, y compris la liste des \omterm{def:nw:where-the-crypto-venues-live:as}{systèmes autonomes} que traverse une route, et choisit parmi les routes selon une politique, non selon la latence. Une
\emph{dorsale cloud}\index{dorsale cloud} est le réseau privé d'un fournisseur entre ses régions, qui achemine le trafic entre elles (peering,
réplication, ses propres services) selon le routage du fournisseur.
\end{definition}

Trois types de chemins relient deux régions. L'internet public, où l'opérateur de chaque saut choisit le suivant selon sa politique commerciale et où le BGP
transporte l'accessibilité, non le délai~: les chemins courts sont un effet de bord. La dorsale du fournisseur, conçue pour la capacité et la résilience, et
dans le blog d'AWS, non pour la latence la plus faible. Et les lignes privées~: \omterm{def:nw:long-haul-fibre:dark}{fibre noire} ou longueurs d'onde (chapitre~13) et radio (chapitre~14), achetées
directement ou auprès d'un réseau spécialisé dans la faible latence qui se relie aux régions du fournisseur par son produit d'\omterm{def:nw:colocation-products-and-how-they-are-sold:mmr}{interconnexion}, comme le fait le spécialiste du blog.

\begin{dated}{2026-09}{Routes entre régions et règles de doublons, d'après des sources publiques}\label{dat:nw:cross-region-crypto-networks:routes}
\begin{itemize}
\item AWS (juin 2026)~: médianes des allers-retours TCP depuis Tokyo le 28 avril 2026, par peering VPC et par Direct Connect avec Avelacom~: Singapour 68.0 et
  65.4~ms~; Virginie 150.2 et 135.4~; Francfort 224.3 et 136.5~; Londres 210.0 et 139.6~; Stockholm 245.3 et 124.3.
\item API spot de Binance~: un identifiant d'ordre client est unique parmi les ordres ouverts~; un ordre qui en réutilise un ne peut être accepté qu'une fois
  l'ordre précédent exécuté, sinon il est rejeté («~Duplicate order sent.~»).
\end{itemize}
\end{dated}

\section{Hébergeurs spécialisés à faible latence et serveurs virtuels}

Entre Tokyo et l'Europe, le spécialiste du blog fait gagner 70 à 121~millisecondes d'aller-retour par rapport à la dorsale. Pour une stratégie qui négocie entre une
plateforme de Tokyo et une plateforme européenne, c'est la différence entre voir les cours de l'autre marché avec un dixième de seconde de retard et les voir à
temps pour compter. Le spécialiste est un \omterm{def:nw:buying-connectivity:extranet}{extranet financier} (chapitre~15) pour les régions cloud~: des lignes privées entre les hubs, livrées dans les
régions du fournisseur par son produit d'\omterm{def:nw:colocation-products-and-how-they-are-sold:mmr}{interconnexion}.

La \omterm{def:nw:where-the-crypto-venues-live:lottery}{loterie des instances} du chapitre~17 s'applique aux deux bouts~: la route inter-régions la plus rapide est gâchée sur une machine située loin de la plateforme
dans sa propre région.

\section{Mise en concurrence du même message sur plusieurs routes}

\begin{definition}[Mise en concurrence des routes, clé d'idempotence]\label{def:nw:cross-region-crypto-networks:idem}
La \emph{mise en concurrence des routes}\index{mise en concurrence des routes} (route racing) envoie le même message sur plusieurs routes à la fois et retient la
première copie arrivée. Une \emph{clé d'idempotence}\index{clé d'idempotence} est un identifiant que l'émetteur attache à chaque copie d'un même message (pour
les ordres, l'identifiant d'ordre client) afin que le récepteur puisse reconnaître les copies et n'en traiter qu'une.
\end{definition}

\begin{proposition}[Ce que la mise en concurrence rapporte]\label{prop:nw:cross-region-crypto-networks:race}
La première de $k$ copies arrive au minimum des latences des routes. Si les délais des routes étaient parfaitement corrélés, ce minimum serait le
délai de la route la plus rapide et la mise en concurrence ne rapporterait rien~; plus elles sont décorrélées, plus la queue du minimum passe sous la queue de chaque route
prise isolément, puisque le minimum ne dépasse un niveau que si toutes les routes le dépassent.
\end{proposition}

\begin{proof}
$P(\min_i X_i > t) = P(X_1 > t, \dots, X_k > t)$, qui égale la plus petite probabilité de survie isolée lorsque les routes évoluent ensemble et tombe
au produit de ces probabilités lorsqu'elles sont indépendantes.
\end{proof}

Le modèle met en concurrence trois routes unidirectionnelles de Tokyo à Singapour~: celle du spécialiste et celle du backbone, à la moitié de leurs allers-retours publiés (32.7 et
\qty{34.0}{\milli\second}), chacune avec un 99\up{e}~centile supposé supérieur de 15~\% à sa médiane, et une route par internet à \qty{35}{\milli\second} avec une
queue plus épaisse (une hypothèse). Une copule gaussienne les relie par la corrélation $\rho$, qui représente une cause commune comme un câble partagé ou une congestion.

\begin{omfigure}
\begin{tikzpicture}
\begin{axis}[width=11cm,height=5.8cm,xmin=0,xmax=0.9,ymin=0,ymax=2300,xlabel={corrélation des délais des routes ($\rho$)},
  ylabel={p99 économisé (\si{\micro\second})},tick label style={font=\footnotesize,/pgf/number format/1000 sep={\,}},label style={font=\small},grid=major,grid style={gray!25},
  legend pos=north east,legend style={font=\footnotesize},table/col sep=comma]
\addplot[omDef,thick,mark=*,mark size=1.4pt] table[x=rho,y=two]{figdata/networks/19-cross-region-crypto-networks/race.csv};
\addplot[omThm,thick,mark=square*,mark size=1.4pt] table[x=rho,y=three]{figdata/networks/19-cross-region-crypto-networks/race.csv};
\legend{deux routes (spécialiste et backbone),trois routes (avec internet)}
\end{axis}
\end{tikzpicture}

\omcaption{Simulation~: latence unidirectionnelle au 99\up{e}~centile économisée par la mise en concurrence, de Tokyo à Singapour, par rapport à la meilleure route isolée, à mesure que la corrélation des
routes augmente. À $\rho = 0$, deux routes économisent 1.6 et trois \qty{2.1}{\milli\second}~; à $\rho = 0.9$, les deux économisent environ
\qty{50}{\micro\second}. Données~: \texttt{fig\_race.py}, \texttt{nw\_race.race\_rows()}.}\label{fig:nw:cross-region-crypto-networks:race}
\end{omfigure}

\omcode{code/firm/routerace/firm_routerace.py}{35}{50}{Routes corrélées par une copule gaussienne, première arrivée, et centile économisé
par rapport à la meilleure route isolée.}\label{lst:nw:cross-region-crypto-networks:race}

La mise en concurrence rapporte dans la queue, et seulement tant que les routes échouent de façon indépendante. Avec des routes indépendantes, la deuxième route fait gagner \qty{1623}{\micro\second} au
99\up{e}~centile et la troisième encore \qty{518}{\micro\second}~; à $\rho = 0.5$, 771 et 90~; à 0.9, 53 et rien. Pour un coût supposé de 10~USD par
mois et par microseconde de 99\up{e}~centile, la troisième route vaut 5\,182~USD par mois si les routes sont indépendantes, 902 à $\rho = 0.5$ et
rien à 0.9~: c'est le prix au-delà duquel la mise en concurrence ne rapporte plus.

Ce qui arrive en second pose problème. La plateforme reçoit plusieurs fois le même ordre~; la \omterm{def:nw:cross-region-crypto-networks:idem}{clé d'idempotence} doit faire qu'il ne soit traité qu'une fois. L'interface spot de Binance
est typique de la règle à vérifier~: un identifiant d'ordre client doit être unique parmi les ordres ouverts, et une copie n'est rejetée comme doublon que
tant que le premier ordre est ouvert. Une fois la première copie exécutée, une copie tardive portant le même identifiant est un nouvel ordre.

\begin{omfigure}
\begin{tikzpicture}
\begin{axis}[width=11cm,height=5.4cm,xmin=0,xmax=8000,ymin=0,ymax=2.1,xlabel={temps pour que la première copie soit exécutée (\si{\micro\second})},x tick label style={/pgf/number format/1000 sep={\,}},
  ylabel={doublons acceptés par ordre},tick label style={font=\footnotesize},label style={font=\small},grid=major,grid style={gray!25},
  table/col sep=comma,scaled x ticks=false]
\addplot[omMeth,thick,mark=*,mark size=1.4pt] table[x=fill_us,y=accepted]{figdata/networks/19-cross-region-crypto-networks/dups.csv};
\end{axis}
\end{tikzpicture}

\omcaption{Simulation~: avec trois copies mises en concurrence ($\rho = 0.5$) et une plateforme qui ne rejette un identifiant en double que tant que le premier ordre est ouvert, le
nombre moyen de copies tardives acceptées comme nouveaux ordres, en fonction du temps que met la première copie à être exécutée. Un ordre exécuté dès son arrivée est
exécuté trois fois. Données~: \texttt{fig\_race.py}, \texttt{nw\_race.duplicate\_rows()}.}\label{fig:nw:cross-region-crypto-networks:dups}
\end{omfigure}

\begin{proposition}[Mise en concurrence d'ordres exécutables]\label{prop:nw:cross-region-crypto-networks:dup}
Sous une règle qui ne rejette un identifiant répété que tant que le premier ordre est ouvert, un ordre mis en concurrence et exécuté dans un délai $f$ après
l'arrivée de la première copie est exécuté à nouveau par chaque copie postérieure arrivant après $f$. Pour les ordres exécutés dès leur arrivée ($f = 0$), $k$ copies mises en concurrence font $k$ ordres.
\end{proposition}

\begin{proof}
Une copie tardive qui trouve le premier ordre déjà exécuté ne trouve aucun ordre ouvert portant son identifiant et est acceptée. Avec $f = 0$, toute copie ultérieure est tardive.
\end{proof}

Le modèle rend le danger concret~: trois copies d'un ordre exécuté dès son arrivée donnent trois exécutions~; un ordre qui met une milliseconde à être exécuté produit
encore 1.57 exécutions indésirables en moyenne~; huit millisecondes, 0.07. La mise en concurrence est donc sûre pour les messages idempotents
sur la plateforme (une annulation, une requête, des données de marché) ou sous une règle qui se souvient des identifiants après la fin de l'ordre, et dangereuse pour les ordres d'exécution sous une
règle limitée aux ordres ouverts. Le traitement idempotent du livre~13 est ce que la plateforme doit faire~; la leçon de ce chapitre est de lire la règle de la plateforme avant de mettre en concurrence.

\begin{omfigure}
\begin{tikzpicture}[font=\footnotesize,b/.style={draw,rounded corners=2pt,minimum height=0.55cm,align=center,inner sep=3pt},>=Stealth]
\node[b,fill=omDef!15] (s) at (0,0) {firme (Tokyo)\\ ordre, clé K};
\node[b,fill=omThm!15,minimum height=1.6cm] (v) at (9.6,0) {plateforme\\ (Singapour)};
\draw[->,thick] (1.2,0.3) -- node[above,font=\scriptsize]{spécialiste~: premier, accepté} (8.7,0.5);
\draw[->,thick,gray] (1.2,0) -- node[above,font=\scriptsize]{backbone~: deuxième} (8.7,0);
\draw[->,thick,gray,dashed] (1.2,-0.3) -- node[below,font=\scriptsize]{internet~: troisième} (8.7,-0.5);
\node[font=\scriptsize,align=left,anchor=west] at (10.6,0.3) {K ouvert~: rejeter les copies};
\node[font=\scriptsize,align=left,anchor=west,omMeth] at (10.6,-0.3) {K exécuté~: nouvel ordre accepté};
\end{tikzpicture}

\omcaption{\omterm{def:nw:cross-region-crypto-networks:idem}{Mise en concurrence des routes}, schématiquement~: trois copies d'un ordre avec une même \omterm{def:nw:cross-region-crypto-networks:idem}{clé d'idempotence}~; la plateforme garde la première et doit reconnaître les autres.
Sous une règle limitée aux ordres ouverts, elle les rejette tant que l'ordre est ouvert et les accepte comme nouveaux ordres une fois qu'il a été exécuté.}\label{fig:nw:cross-region-crypto-networks:scheme}
\end{omfigure}

\begin{method}[Mise en concurrence des routes entre régions]\label{met:nw:cross-region-crypto-networks:race}
\begin{enumerate}
\item Mesurer la distribution de latence de chaque route candidate, et pas seulement sa médiane, ainsi que la corrélation de leurs délais (partagent-elles un câble, un
  point d'atterrissement, un point de peering~?).
\item Ne mettre en concurrence que les routes dont les queues ne sont pas trop corrélées, et chiffrer chaque route supplémentaire par rapport à la queue qu'elle économise.
\item Lire la règle de la plateforme sur les doublons. Mettre en concurrence les annulations, les modifications et les requêtes sans restriction~; ne mettre en concurrence les ordres d'exécution que si la plateforme se souvient des identifiants après
  la clôture de l'ordre, ou si la logique propre de la firme tolère un doublon.
\item Journaliser l'arrivée de chaque copie (à partir des accusés de réception de la plateforme) et remesurer la corrélation des routes~; elle change lorsqu'un câble est coupé ou
  réparé.
\end{enumerate}
\end{method}

\section{Tutoriel~: hubs et mises en concurrence}\label{tut:nw:cross-region-crypto-networks:tut}

\begin{tutorial}
\textbf{Objectif.} Mesurer l'\omterm{def:nw:cross-region-crypto-networks:infl}{allongement du chemin} entre les hubs et simuler la \omterm{def:nw:cross-region-crypto-networks:idem}{mise en concurrence des routes} avec ses doublons.
\textbf{État final~:} \cref{fig:nw:cross-region-crypto-networks:hubs,fig:nw:cross-region-crypto-networks:race,fig:nw:cross-region-crypto-networks:dups}
et \cref{tab:nw:cross-region-crypto-networks:hubs}.
\begin{tutsteps}
\item \textbf{Hubs.} Les villes des régions sont les lignes de \texttt{data/networks/sites.csv}~; \texttt{nw\_race.hub\_rows()} calcule les \omterm{def:nw:the-north-american-map:geodesic}{distances géodésiques}, les planchers et
  l'\omterm{def:nw:cross-region-crypto-networks:infl}{allongement du chemin} par rapport aux allers-retours publiés.
\item \textbf{Routes.} \texttt{ROUTES} contient trois routes à partir des médianes publiées et des queues supposées~; \texttt{first\_arrival}
  (\cref{lst:nw:cross-region-crypto-networks:race}) les tire de façon corrélée.
\item \textbf{Gain et prix.} \texttt{race\_rows()} et \texttt{break\_even()}.
\item \textbf{Doublons.} \texttt{duplicates} applique la règle limitée aux ordres ouverts~; \texttt{duplicate\_rows()} balaie le temps d'exécution.
\end{tutsteps}
\textbf{Pour aller plus loin.} Donner à la route du backbone vers Francfort ses 224~ms publiés et la mettre en concurrence avec celle du spécialiste~; fixer la règle à
«~toujours~» et constater que les doublons disparaissent.
\end{tutorial}

\section{Construction~: le moteur de mise en concurrence des routes}\label{bld:nw:cross-region-crypto-networks:routerace}

\begin{build}
\bfield{Rôle} La latence, le coût et le risque de l'envoi de messages sur plusieurs routes entre régions~: l'entrée des lignes inter-régions
du plan du chapitre~29.
\bfield{Interface} \texttt{firm\_routerace}~: \texttt{Route}, \texttt{first\_arrival}, \texttt{percentile\_gain}, \texttt{break\_even\_price},
\texttt{duplicates}, \texttt{path\_inflation}.
\bfield{Règles} La corrélation est explicite~; les tirages de routes utilisent des nombres aléatoires communs d'un nombre de routes à l'autre~; le traitement des doublons suit une règle de plateforme énoncée~;
les simulations le disent.
\bfield{Tests d'acceptation} \texttt{code/firm/routerace/tests/}~: la médiane et le 99\up{e}~centile des lois marginales, la première arrivée, l'absence de gain en
corrélation parfaite, les règles de doublons vérifiées à la main, l'\omterm{def:nw:cross-region-crypto-networks:infl}{allongement du chemin} d'une route parfaite.
\bfield{Pour aller plus loin} Ajuster les délais de route et leur corrélation à partir des accusés de réception journalisés~; un émetteur qui met ses routes en concurrence annule et n'envoie chaque ordre qu'une fois.
\end{build}

\begin{omsources}
\item Blog AWS for Industries, «~Ultra-low-latency cross-Region crypto trading with Avelacom and AWS~» (juin 2026)~; documentation des AWS Regions.
\item RFC 4271 (BGP-4)~; Binance spot API documentation (identifiants d'ordres client et ordres en double).
\end{omsources}

\section{Exercices}

\begin{exercise}[$\star$]\label{exo:nw:cross-region-crypto-networks:1}
Quel est l'\omterm{def:nw:cross-region-crypto-networks:infl}{allongement du chemin} du backbone et celui du spécialiste de Tokyo à Francfort~?
\end{exercise}

\begin{exercise}[$\star$]\label{exo:nw:cross-region-crypto-networks:2}
Pourquoi BGP ne choisit-il pas le chemin le plus rapide~?
\end{exercise}

\begin{exercise}[$\star$]\label{exo:nw:cross-region-crypto-networks:3}
Quel aller-retour le spécialiste permet-il d'économiser par rapport au backbone, de Tokyo à chacun des hubs européens~?
\end{exercise}

\begin{exercise}[$\star\star$]\label{exo:nw:cross-region-crypto-networks:4}
Deux routes partagent un câble sous-marin sur la moitié de leur longueur. Que cela change-t-il à la valeur de leur mise en concurrence~?
\end{exercise}

\begin{exercise}[$\star\star$]\label{exo:nw:cross-region-crypto-networks:5}
Quels messages peut-on mettre en concurrence en toute sécurité sous une règle de doublons limitée aux ordres ouverts, et lesquels ne le peut-on pas~?
\end{exercise}

\begin{exercise}[$\star\star$]\label{exo:nw:cross-region-crypto-networks:6}
Une stratégie négocie Binance à Tokyo contre une plateforme européenne. Où son moteur doit-il tourner~?
\end{exercise}

\begin{exercise}[$\star\star\star$]\label{exo:nw:cross-region-crypto-networks:7}
\emph{Programmation.} Avec \texttt{firm.routerace}, trouver la corrélation au-delà de laquelle la troisième route fait gagner moins de \qty{100}{\micro\second} au 99e
centile.
\end{exercise}

\begin{exercise}[$\star\star\star$]\label{exo:nw:cross-region-crypto-networks:8}
\emph{Trouver la faille.} «~Nous mettons chaque ordre en concurrence sur trois routes~; la plateforme rejette les doublons, donc c'est une latence gratuite.~»
\end{exercise}

\section{Problème~: trois routes vers Singapour}

\begin{problem}[{Problème du week-end --- ce que rapporte la mise en concurrence}]\label{pb:nw:cross-region-crypto-networks:1}
Une société de Tokyo envoie des ordres à une plateforme de Singapour par trois routes au plus~: un réseau spécialisé (\qty{32.7}{\milli\second} à l'aller, médiane), le
backbone du fournisseur (\qty{34.0}{\milli\second}) et l'internet (\qty{35}{\milli\second})~; les 99\up{e}~centiles des deux routes privées sont supérieurs de 15~\%
à leurs médianes, celui de l'internet de 30~\%. Une microseconde de 99\up{e}~centile lui vaut 10~USD par mois.

\medskip\noindent\textbf{Partie I --- Les planchers.}
\begin{enumerate}
\item Quels sont les planchers géodésique, du vide et de l'aller-retour en fibre entre Tokyo et Singapour~?
\item Quel est l'\omterm{def:nw:cross-region-crypto-networks:infl}{allongement du chemin} de chaque route publiée~?
\item Quelle part de l'aller-retour du spécialiste dépasse le plancher de la fibre~?
\item Pourquoi la médiane de la route internet est-elle une hypothèse plutôt qu'un chiffre publié~?
\end{enumerate}

\medskip\noindent\textbf{Partie II --- La mise en concurrence.}
\begin{enumerate}[resume]
\item Que fait gagner la mise en concurrence de deux routes au 99\up{e}~centile, pour $\rho = 0$, 0.5 et 0.9~?
\item Et trois routes~?
\item Que vaut la troisième route par mois, à chaque corrélation~?
\item Que fait la mise en concurrence à la médiane~?
\end{enumerate}

\medskip\noindent\textbf{Partie III --- Les doublons.}
\begin{enumerate}[resume]
\item Sous une règle limitée aux ordres ouverts, combien de copies tardives d'un ordre qui s'exécute à l'arrivée sont-elles exécutées~?
\item Et pour un ordre qui met une milliseconde, quatre ou huit à s'exécuter~?
\item Lesquels des messages de la société faudrait-il mettre en concurrence~?
\item Quelle règle faudrait-il demander à la plateforme~?
\end{enumerate}

\medskip\noindent\textbf{Partie IV --- Le verdict.}
\begin{enumerate}[resume]
\item Énoncer le \emph{résultat nommé}~: le gain au 99\up{e}~centile de la mise en concurrence de deux ou trois routes de corrélation $\rho$, et le prix par route au-delà
  duquel la mise en concurrence ne rapporte pas.
\item Quelle donnée fait le plus bouger la réponse~?
\item Comment estimer $\rho$~?
\item Que devient $\rho$ lorsqu'un câble est coupé~?
\item Faudrait-il mettre en concurrence une route vers Francfort, où le backbone est plus lent de 88~ms que le spécialiste~?
\item Que fait la position de la société dans chaque région (chapitre~17) à tout cela~?
\item Où cela s'insère-t-il dans le plan du chapitre~29~?
\item En une phrase~: quand la mise en concurrence rapporte-t-elle~?
\end{enumerate}
\end{problem}

\section{Questions d'entretien}

\begin{interviewq}[$\star$ \iqroles{développeur}]\label{iq:nw:cross-region-crypto-networks:1}
Pourquoi l'aller-retour de Tokyo à Francfort sur la dorsale du cloud dépasse-t-il plus du double du plancher physique~?
\end{interviewq}

\begin{interviewq}[$\star\star$ \iqroles{développeur}]\label{iq:nw:cross-region-crypto-networks:2}
Quand paieriez-vous un réseau spécialisé entre deux régions cloud~?
\end{interviewq}

\begin{interviewq}[$\star\star$ \iqroles{développeur, chercheur}]\label{iq:nw:cross-region-crypto-networks:3}
Vous envoyez le même message sur deux routes. Qu'est-ce que cela vous apporte, et de quoi cela dépend-il~?
\end{interviewq}

\begin{interviewq}[$\star\star$ \iqroles{développeur}]\label{iq:nw:cross-region-crypto-networks:4}
Qu'est-ce qu'une \omterm{def:nw:cross-region-crypto-networks:idem}{clé d'idempotence}, et comment peut-elle ne pas vous protéger~?
\end{interviewq}

\begin{interviewq}[$\star\star$ \iqroles{trader}]\label{iq:nw:cross-region-crypto-networks:5}
Comment une centaine de millisecondes entre Tokyo et l'Europe façonne-t-elle une stratégie crypto entre plateformes~?
\end{interviewq}

\begin{interviewq}[$\star\star\star$ \iqroles{développeur, chercheur}]\label{iq:nw:cross-region-crypto-networks:6}
Concevez une couche d'envoi d'ordres qui met les routes en concurrence sans jamais exécuter un ordre deux fois.
\end{interviewq}
